\chapter[Worked Examples]{Worked Examples}\label{ch:sec:examples}
This section demonstrates the complete mathematical pipeline on three
finite configurations.  The first exhibits the ideal case in detail;
the second and third illustrate failure and progressive restoration
more concisely.  All calculations are reproducible by hand.
\section{The cycle factorization criterion}
On a tree, Theorem~\ref{ch:thm:cycle-factorization} has no cycle equations,
so every positive edge labelling factors.  After choosing one positive
weight at a root, the remaining weights are obtained uniquely by
propagating the equations $\beta(c)=n_{mc}w(m)$ along the unique paths.
For a four-cycle
$m_0-c_1-m_1-c_2-m_0$, the criterion reduces to $n_{m_0c_1}n_{m_1c_2}=n_{m_1c_1}n_{m_0c_2}.$ The labels $(2,3,6,4)$ in the cyclic edge order
$(m_0c_1,m_1c_1,m_1c_2,m_0c_2)$ are compatible because
$2\cdot6=3\cdot4$.  For example, taking $w(m_0)=1$ gives
$\beta(c_1)=2$, $w(m_1)=2/3$, and $\beta(c_2)=4$.
Replacing $6$ by $5$ makes the products $10$ and $12$; no positive
vertex weights can then satisfy all four edge equations.
\section{Example 1: A universally optimal encoder}
Let $\Mset=\{m_1,m_2,m_3\}$ and $\Cset=\{c_1,c_2\}$ with incidence
$I=\{(m_1,c_1),(m_1,c_2),(m_2,c_1),(m_3,c_2)\}$, so
$\cloak(c_1)=\{m_1,m_2\}$, $\cloak(c_2)=\{m_1,m_3\}$.  Choose
$w=(1/2,1/3,1/6)$, $\beta=(4,2)$.  The induced multiplicity matrix is
\[
N=\begin{pmatrix}8&4\\12&0\\0&12\end{pmatrix},
\qquad |\Kset|=12.
\]
Since $w$ sums to $1$, $\pi=w$.  Universal optimality follows from
$\pi(m)n_{mc}=4$ on $\cloak(c_1)$ and $2$ on $\cloak(c_2)$.
Logarithmic potentials are $u(m)=\log w(m)$, $v(c)=\log\beta(c)$.
With $V_I=\mathbb R^4$ ordered by $(m_1c_1,m_1c_2,m_2c_1,m_3c_2)$,
\[
\omega \approx (2.0794,1.3863,2.4849,2.4849),\qquad
Y = A(\log\pi)+\omega = (\log4,\log2,\log4,\log2).
\]
Since $Y$ has the form $(a,b,a,b)$, $Y\in\operatorname{Im}(C)$ and
$\omega\in\operatorname{Im}(\Phi)$.  Hence $\min E_\omega=0$,
$E_{\mathrm{red}}(u)=0$ at $u=-\log\pi$, $QY=0$, $q(Y)=0$, and
$\delta(Y)=0$.
This example confirms mutual consistency of the WIS representation, the
canonical linear representation, the least-squares formulation,
convexity, the quotient geometry, and the defect in the ideal case.
\section{Example 2: Breaking universal optimality}
Modify a single multiplicity in Example~1: set $n_{m_1c_1}=6$,
$n_{m_1c_2}=6$, preserving row sums:
\[
N'=\begin{pmatrix}\mathbf6&\mathbf6\\12&0\\0&12\end{pmatrix}.
\]
With the same prior, $\pi n$ is not constant on cloaks
($3\neq4$, $3\neq2$), so the system is not universally optimal.
\[
\omega' = (\log6,\log6,\log12,\log12),\qquad
Y' \approx (1.0986,1.0986,1.3863,0.6931),
\]
We have $Y'\notin\operatorname{Im}(C)$.  However, the incidence graph
is a tree, so Theorem~\ref{ch:thm:cycle-factorization}
implies $\omega'\in\operatorname{Im}(\Phi)$ and hence
$\min E_{\omega'}=0$: the modified multiplicities still admit an exact
WIS factorization, but its message weights are not proportional to the
fixed prior.  The prior-dependent defect detects precisely this failure:
$\delta(Y') = \|QY'\| \approx 0.3515 > 0.$
All structures (incidence spaces, lifting operators, $Q$, quotient)
remain well defined; only the gauge class is non-zero.  By
Theorem~\ref{ch:thm:stability}, any perturbation of $\omega$ produces a
proportionally bounded change in the reconstructed potentials.
\section{Example 3: Progressive restoration}
Starting from Example~2, correct $n_{m_1c_1}$ toward the WIS target
$8$ (with $n_{m_1c_2}=12-n_{m_1c_1}$):
\[
\begin{array}{cccc}
\text{Step} & n_{m_1c_1} & \delta(Y) & \delta(Y)^2 \\\hline
\text{Initial} & 6 & 0.3515 & 0.1235 \\
\text{Step 1} & 7 & 0.1839 & 0.0338 \\
\text{Final} & 8 & 0 & 0
\end{array}
\]
The defect decreases monotonically to zero as the WIS is restored.
The reduced energy $E_{\mathrm{red}}$ remains convex throughout
(Theorem~\ref{ch:thm:convexity}), so the unique minimizer moves
continuously toward the WIS-compatible potentials.
\paragraph{Stability constant.}
We compute $C=\|(\Phi^{\mathsf T}\Phi)^\dagger\Phi^{\mathsf T}\|$ for
this configuration.  With $|M|=3$, $|C|=2$, $|I|=4$, the operator
$\Phi^{\mathsf T}\Phi$ is $5\times5$.  Its positive eigenvalues are
approximately $0.3820$, $1.3820$, $2.6180$, and $3.6180$.
Equivalently, the smallest nonzero singular value of $\Phi$ is
$\sigma_{\min}^+\approx0.6180$.  Since
$(\Phi^{\mathsf T}\Phi)^\dagger\Phi^{\mathsf T}=\Phi^\dagger$, the
exact operator norm relevant to Theorem~\ref{ch:thm:stability} is
\[
C=\|\Phi^\dagger\|=\frac{1}{\sigma_{\min}^+}
\approx 1.6180.
\]
Thus perturbations of $\omega$ are amplified by at most this constant
in the minimum-norm reconstruction.
\section{Discussion}
The three examples illustrate the complete pipeline:
\[
\text{Encoder}\;\to\; N \;\to\; \text{WIS}
\;\to\; \omega \;\to\; \Phi
\;\to\; \text{least-squares}
\;\to\; q(Y) \;\to\; \delta(Y).
\]
In Example~1 every step is consistent and $\delta=0$.  Example~2 shows
that exact WIS factorization may coexist with a positive
prior-dependent universal optimality defect.  Example~3 confirms
monotonic decrease of that defect and stability of the reconstruction.
The defect therefore provides a quantitative measure of distance from
universal optimality, computed from the incidence structure,
multiplicities, and prior.
\endinput
